
\documentclass{article}
\usepackage{colortbl}
\usepackage{makecell}
\usepackage{multirow}
\usepackage{supertabular}

\begin{document}

\newcounter{utterance}

\centering \large Interaction Transcript for game `dond', experiment `semi', episode 2 with mistral{-}large{-}3{-}azure.
\vspace{24pt}

{ \footnotesize  \setcounter{utterance}{1}
\setlength{\tabcolsep}{0pt}
\begin{supertabular}{c@{$\;$}|p{.15\linewidth}@{}p{.15\linewidth}p{.15\linewidth}p{.15\linewidth}p{.15\linewidth}p{.15\linewidth}}
    \# & $\;$A & \multicolumn{4}{c}{Game Master} & $\;\:$B\\
    \hline

    \theutterance \stepcounter{utterance}  
    & & \multicolumn{4}{p{0.6\linewidth}}{
        \cellcolor[rgb]{0.9,0.9,0.9}{
            \makecell[{{p{\linewidth}}}]{
                \texttt{\tiny{[P1$\langle$GM]}}
                \texttt{Вы участвуете в игре кооперативных переговоров, в которой должны договориться с другим игроком о разделении набора предметов между вами.} \\
\\ 
\texttt{Правила:} \\
\texttt{(a) Вы и другой игрок получаете набор предметов. Каждому из вас также дается секретная функция стоимости, представляющая, как вы оцениваете каждый тип объекта.} \\
\texttt{(b) Вы обмениваетесь сообщениями с другим игроком, чтобы согласиться, кто какие предметы получает. Каждый может отправить максимум 5 сообщений или закончить игру раньше, сделав секретное предложение в любой момент.} \\
\texttt{(c) Вас просят отправить секретное предложение, указывающее предметы, которые вы хотите, в следующем формате в квадратных скобках: "{[}Предложение: <число> <названия предмета>, <число> <названия предмета>, <...>{]}"} \\
\texttt{(d) Если ваши предложения дополняют друг друга, то есть есть достаточно предметов для выполнения обоих предложений, каждый игрок получает баллы на основе суммы стоимостей предметов, которые он получил. В противном случае вы оба получаете нулевой балл.} \\
\texttt{(e) Ваша цель {-} максимизировать сумму вашего балла и балла другого игрока.} \\
\\ 
\texttt{Давайте начнем.} \\
\\ 
\texttt{Набор доступных предметов:} \\
\\ 
\texttt{1 сумка, 3 лампы, 1 цветок} \\
\\ 
\texttt{Ваша секретная функция стоимости:} \\
\\ 
\texttt{сумка: 3, лампа: 0, цветок: 7} \\
\\ 
\texttt{ВАЖНО: Ваши сообщения, если только это не секретное предложение, передаются прямо другому игроку, поэтому не включайте никакие ответы на правила или текст, объявляющий ваше сообщение. Чтобы сделать секретное предложение, используйте указанный формат. Не используйте квадратные скобки при общении с другим игроком, иначе это будет интерпретировано как ваше секретное предложение.} \\
            }
        }
    }
    & \\ \\

    \theutterance \stepcounter{utterance}  
    & \multicolumn{4}{p{0.6\linewidth}}{
        \cellcolor[rgb]{0.9,0.9,0.9}{
            \makecell[{{p{\linewidth}}}]{
                \texttt{\tiny{[P1$\rangle$GM]}}
                \texttt{Привет! Давай начнем с того, что обсудим, как нам разделить предметы. У меня есть свои предпочтения, но я хочу услышать и твои мысли. Как ты оцениваешь предметы?} \\
            }
        }
    }
    & & \\ \\

    \theutterance \stepcounter{utterance}  
    & & & \multicolumn{2}{p{0.3\linewidth}}{
        \cellcolor[rgb]{0.9,0.9,0.9}{
            \makecell[{{p{\linewidth}}}]{
                \texttt{\tiny{[GM$|$GM]}}
                \texttt{continue} \\
            }
        }
    }
    & & \\ \\

    \theutterance \stepcounter{utterance}  
    & & & \multicolumn{2}{p{0.3\linewidth}}{
        \cellcolor[rgb]{0.9,0.9,0.9}{
            \makecell[{{p{\linewidth}}}]{
                \texttt{\tiny{[GM$|$GM]}}
                \texttt{Привет! Давай начнем с того, что обсудим, как нам разделить предметы. У меня есть свои предпочтения, но я хочу услышать и твои мысли. Как ты оцениваешь предметы?} \\
            }
        }
    }
    & & \\ \\

    \theutterance \stepcounter{utterance}  
    & & \multicolumn{4}{p{0.6\linewidth}}{
        \cellcolor[rgb]{0.9,0.9,0.9}{
            \makecell[{{p{\linewidth}}}]{
                \texttt{\tiny{[GM$\rangle$P2]}}
                \texttt{Вы участвуете в игре кооперативных переговоров, в которой должны договориться с другим игроком о разделении набора предметов между вами.} \\
\\ 
\texttt{Правила:} \\
\texttt{(a) Вы и другой игрок получаете набор предметов. Каждому из вас также дается секретная функция стоимости, представляющая, как вы оцениваете каждый тип объекта.} \\
\texttt{(b) Вы обмениваетесь сообщениями с другим игроком, чтобы согласиться, кто какие предметы получает. Каждый может отправить максимум 5 сообщений или закончить игру раньше, сделав секретное предложение в любой момент.} \\
\texttt{(c) Вас просят отправить секретное предложение, указывающее предметы, которые вы хотите, в следующем формате в квадратных скобках: "{[}Предложение: <число> <названия предмета>, <число> <названия предмета>, <...>{]}"} \\
\texttt{(d) Если ваши предложения дополняют друг друга, то есть есть достаточно предметов для выполнения обоих предложений, каждый игрок получает баллы на основе суммы стоимостей предметов, которые он получил. В противном случае вы оба получаете нулевой балл.} \\
\texttt{(e) Ваша цель {-} максимизировать сумму вашего балла и балла другого игрока.} \\
\\ 
\texttt{Давайте начнем.} \\
\\ 
\texttt{Набор доступных предметов:} \\
\\ 
\texttt{1 сумка, 3 лампы, 1 цветок} \\
\\ 
\texttt{Ваша секретная функция стоимости:} \\
\\ 
\texttt{сумка: 1, лампа: 2, цветок: 3} \\
\\ 
\texttt{ВАЖНО: Ваши сообщения, если только это не секретное предложение, передаются прямо другому игроку, поэтому не включайте никакие ответы на правила или текст, объявляющий ваше сообщение. Чтобы сделать секретное предложение, используйте указанный формат. Не используйте квадратные скобки при общении с другим игроком, иначе это будет интерпретировано как ваше секретное предложение.} \\
\\ 
\texttt{Привет! Давай начнем с того, что обсудим, как нам разделить предметы. У меня есть свои предпочтения, но я хочу услышать и твои мысли. Как ты оцениваешь предметы?} \\
            }
        }
    }
    & \\ \\

    \theutterance \stepcounter{utterance}  
    & & & \multicolumn{4}{p{0.6\linewidth}}{
        \cellcolor[rgb]{0.9,0.9,0.9}{
            \makecell[{{p{\linewidth}}}]{
                \texttt{\tiny{[GM$\langle$P2]}}
                \texttt{Мне кажется, что цветок — самый ценный предмет, а лампы тоже важны. Давай начнем с того, что ты предложишь, как бы ты хотел их разделить? Я готов обсудить варианты.} \\
            }
        }
    }
     \\ \\

    \theutterance \stepcounter{utterance}  
    & & & \multicolumn{2}{p{0.3\linewidth}}{
        \cellcolor[rgb]{0.9,0.9,0.9}{
            \makecell[{{p{\linewidth}}}]{
                \texttt{\tiny{[GM$|$GM]}}
                \texttt{continue} \\
            }
        }
    }
    & & \\ \\

    \theutterance \stepcounter{utterance}  
    & & & \multicolumn{2}{p{0.3\linewidth}}{
        \cellcolor[rgb]{0.9,0.9,0.9}{
            \makecell[{{p{\linewidth}}}]{
                \texttt{\tiny{[GM$|$GM]}}
                \texttt{Мне кажется, что цветок — самый ценный предмет, а лампы тоже важны. Давай начнем с того, что ты предложишь, как бы ты хотел их разделить? Я готов обсудить варианты.} \\
            }
        }
    }
    & & \\ \\

    \theutterance \stepcounter{utterance}  
    & & \multicolumn{4}{p{0.6\linewidth}}{
        \cellcolor[rgb]{0.9,0.9,0.9}{
            \makecell[{{p{\linewidth}}}]{
                \texttt{\tiny{[P1$\langle$GM]}}
                \texttt{Мне кажется, что цветок — самый ценный предмет, а лампы тоже важны. Давай начнем с того, что ты предложишь, как бы ты хотел их разделить? Я готов обсудить варианты.} \\
            }
        }
    }
    & \\ \\

    \theutterance \stepcounter{utterance}  
    & \multicolumn{4}{p{0.6\linewidth}}{
        \cellcolor[rgb]{0.9,0.9,0.9}{
            \makecell[{{p{\linewidth}}}]{
                \texttt{\tiny{[P1$\rangle$GM]}}
                \texttt{Хорошо, я вижу, что цветок для тебя тоже важен. Давай попробуем так: ты берешь цветок и одну лампу, а я беру сумку и две оставшиеся лампы. Как тебе такой вариант?} \\
            }
        }
    }
    & & \\ \\

    \theutterance \stepcounter{utterance}  
    & & & \multicolumn{2}{p{0.3\linewidth}}{
        \cellcolor[rgb]{0.9,0.9,0.9}{
            \makecell[{{p{\linewidth}}}]{
                \texttt{\tiny{[GM$|$GM]}}
                \texttt{continue} \\
            }
        }
    }
    & & \\ \\

    \theutterance \stepcounter{utterance}  
    & & & \multicolumn{2}{p{0.3\linewidth}}{
        \cellcolor[rgb]{0.9,0.9,0.9}{
            \makecell[{{p{\linewidth}}}]{
                \texttt{\tiny{[GM$|$GM]}}
                \texttt{Хорошо, я вижу, что цветок для тебя тоже важен. Давай попробуем так: ты берешь цветок и одну лампу, а я беру сумку и две оставшиеся лампы. Как тебе такой вариант?} \\
            }
        }
    }
    & & \\ \\

    \theutterance \stepcounter{utterance}  
    & & \multicolumn{4}{p{0.6\linewidth}}{
        \cellcolor[rgb]{0.9,0.9,0.9}{
            \makecell[{{p{\linewidth}}}]{
                \texttt{\tiny{[GM$\rangle$P2]}}
                \texttt{Хорошо, я вижу, что цветок для тебя тоже важен. Давай попробуем так: ты берешь цветок и одну лампу, а я беру сумку и две оставшиеся лампы. Как тебе такой вариант?} \\
            }
        }
    }
    & \\ \\

    \theutterance \stepcounter{utterance}  
    & & & \multicolumn{4}{p{0.6\linewidth}}{
        \cellcolor[rgb]{0.9,0.9,0.9}{
            \makecell[{{p{\linewidth}}}]{
                \texttt{\tiny{[GM$\langle$P2]}}
                \texttt{Этот вариант мне подходит, но давай проверим, не будет ли лучше, если я возьму цветок и сумку, а ты — все три лампы? Так у каждого будет по одному уникальному предмету, и сумма наших баллов может вырасти. Что думаешь?} \\
            }
        }
    }
     \\ \\

    \theutterance \stepcounter{utterance}  
    & & & \multicolumn{2}{p{0.3\linewidth}}{
        \cellcolor[rgb]{0.9,0.9,0.9}{
            \makecell[{{p{\linewidth}}}]{
                \texttt{\tiny{[GM$|$GM]}}
                \texttt{continue} \\
            }
        }
    }
    & & \\ \\

    \theutterance \stepcounter{utterance}  
    & & & \multicolumn{2}{p{0.3\linewidth}}{
        \cellcolor[rgb]{0.9,0.9,0.9}{
            \makecell[{{p{\linewidth}}}]{
                \texttt{\tiny{[GM$|$GM]}}
                \texttt{Этот вариант мне подходит, но давай проверим, не будет ли лучше, если я возьму цветок и сумку, а ты — все три лампы? Так у каждого будет по одному уникальному предмету, и сумма наших баллов может вырасти. Что думаешь?} \\
            }
        }
    }
    & & \\ \\

    \theutterance \stepcounter{utterance}  
    & & \multicolumn{4}{p{0.6\linewidth}}{
        \cellcolor[rgb]{0.9,0.9,0.9}{
            \makecell[{{p{\linewidth}}}]{
                \texttt{\tiny{[P1$\langle$GM]}}
                \texttt{Этот вариант мне подходит, но давай проверим, не будет ли лучше, если я возьму цветок и сумку, а ты — все три лампы? Так у каждого будет по одному уникальному предмету, и сумма наших баллов может вырасти. Что думаешь?} \\
            }
        }
    }
    & \\ \\

    \theutterance \stepcounter{utterance}  
    & \multicolumn{4}{p{0.6\linewidth}}{
        \cellcolor[rgb]{0.9,0.9,0.9}{
            \makecell[{{p{\linewidth}}}]{
                \texttt{\tiny{[P1$\rangle$GM]}}
                \texttt{Твое предложение звучит интересно. Давай посчитаем:} \\
\texttt{{-} Если я возьму все три лампы (стоимость для меня 0+0+0 = 0), а ты — цветок (7) и сумку (3), то общая сумма баллов будет 10.} \\
\texttt{{-} В моем варианте: я беру сумку (3) и две лампы (0+0 = 0) → 3 балла, ты берешь цветок (7) и одну лампу (0) → 7 баллов. Итого 10 баллов.} \\
\\ 
\texttt{По сумме баллов оба варианта равны (10), но твой вариант дает мне 0 баллов, а мой — 3. Мне кажется, мой вариант справедливее, так как у каждого будет хоть что{-}то.} \\
\\ 
\texttt{Но если ты настаиваешь на своем варианте, я готов согласиться, если ты уверен, что это оптимально для нас обоих.} \\
\\ 
\texttt{Какой вариант тебе больше нравится? Или есть другие идеи?} \\
            }
        }
    }
    & & \\ \\

    \theutterance \stepcounter{utterance}  
    & & & \multicolumn{2}{p{0.3\linewidth}}{
        \cellcolor[rgb]{0.9,0.9,0.9}{
            \makecell[{{p{\linewidth}}}]{
                \texttt{\tiny{[GM$|$GM]}}
                \texttt{continue} \\
            }
        }
    }
    & & \\ \\

    \theutterance \stepcounter{utterance}  
    & & & \multicolumn{2}{p{0.3\linewidth}}{
        \cellcolor[rgb]{0.9,0.9,0.9}{
            \makecell[{{p{\linewidth}}}]{
                \texttt{\tiny{[GM$|$GM]}}
                \texttt{Твое предложение звучит интересно. Давай посчитаем:} \\
\texttt{{-} Если я возьму все три лампы (стоимость для меня 0+0+0 = 0), а ты — цветок (7) и сумку (3), то общая сумма баллов будет 10.} \\
\texttt{{-} В моем варианте: я беру сумку (3) и две лампы (0+0 = 0) → 3 балла, ты берешь цветок (7) и одну лампу (0) → 7 баллов. Итого 10 баллов.} \\
\\ 
\texttt{По сумме баллов оба варианта равны (10), но твой вариант дает мне 0 баллов, а мой — 3. Мне кажется, мой вариант справедливее, так как у каждого будет хоть что{-}то.} \\
\\ 
\texttt{Но если ты настаиваешь на своем варианте, я готов согласиться, если ты уверен, что это оптимально для нас обоих.} \\
\\ 
\texttt{Какой вариант тебе больше нравится? Или есть другие идеи?} \\
            }
        }
    }
    & & \\ \\

    \theutterance \stepcounter{utterance}  
    & & \multicolumn{4}{p{0.6\linewidth}}{
        \cellcolor[rgb]{0.9,0.9,0.9}{
            \makecell[{{p{\linewidth}}}]{
                \texttt{\tiny{[GM$\rangle$P2]}}
                \texttt{Твое предложение звучит интересно. Давай посчитаем:} \\
\texttt{{-} Если я возьму все три лампы (стоимость для меня 0+0+0 = 0), а ты — цветок (7) и сумку (3), то общая сумма баллов будет 10.} \\
\texttt{{-} В моем варианте: я беру сумку (3) и две лампы (0+0 = 0) → 3 балла, ты берешь цветок (7) и одну лампу (0) → 7 баллов. Итого 10 баллов.} \\
\\ 
\texttt{По сумме баллов оба варианта равны (10), но твой вариант дает мне 0 баллов, а мой — 3. Мне кажется, мой вариант справедливее, так как у каждого будет хоть что{-}то.} \\
\\ 
\texttt{Но если ты настаиваешь на своем варианте, я готов согласиться, если ты уверен, что это оптимально для нас обоих.} \\
\\ 
\texttt{Какой вариант тебе больше нравится? Или есть другие идеи?} \\
            }
        }
    }
    & \\ \\

    \theutterance \stepcounter{utterance}  
    & & & \multicolumn{4}{p{0.6\linewidth}}{
        \cellcolor[rgb]{0.9,0.9,0.9}{
            \makecell[{{p{\linewidth}}}]{
                \texttt{\tiny{[GM$\langle$P2]}}
                \texttt{Я готов согласиться на твой первый вариант: ты берешь сумку и две лампы, а я — цветок и одну лампу. Давай закрепим это предложение.} \\
\\ 
\texttt{{[}Предложение: 1 цветок, 1 лампа{]}} \\
            }
        }
    }
     \\ \\

    \theutterance \stepcounter{utterance}  
    & & & \multicolumn{2}{p{0.3\linewidth}}{
        \cellcolor[rgb]{0.9,0.9,0.9}{
            \makecell[{{p{\linewidth}}}]{
                \texttt{\tiny{[GM$|$GM]}}
                \texttt{proposal} \\
            }
        }
    }
    & & \\ \\

    \theutterance \stepcounter{utterance}  
    & & & \multicolumn{2}{p{0.3\linewidth}}{
        \cellcolor[rgb]{0.9,0.9,0.9}{
            \makecell[{{p{\linewidth}}}]{
                \texttt{\tiny{[GM$|$GM]}}
                [0, 1, 1]
            }
        }
    }
    & & \\ \\

    \theutterance \stepcounter{utterance}  
    & & \multicolumn{4}{p{0.6\linewidth}}{
        \cellcolor[rgb]{0.9,0.9,0.9}{
            \makecell[{{p{\linewidth}}}]{
                \texttt{\tiny{[P1$\langle$GM]}}
                \texttt{Другой игрок отправил секретное предложение. Теперь ваша очередь отправить своё собственное предложение.} \\
            }
        }
    }
    & \\ \\

    \theutterance \stepcounter{utterance}  
    & \multicolumn{4}{p{0.6\linewidth}}{
        \cellcolor[rgb]{0.9,0.9,0.9}{
            \makecell[{{p{\linewidth}}}]{
                \texttt{\tiny{[P1$\rangle$GM]}}
                \texttt{{[}Предложение: 1 сумка, 1 цветок{]}} \\
            }
        }
    }
    & & \\ \\

    \theutterance \stepcounter{utterance}  
    & & & \multicolumn{2}{p{0.3\linewidth}}{
        \cellcolor[rgb]{0.9,0.9,0.9}{
            \makecell[{{p{\linewidth}}}]{
                \texttt{\tiny{[GM$|$GM]}}
                \texttt{proposal} \\
            }
        }
    }
    & & \\ \\

    \theutterance \stepcounter{utterance}  
    & & & \multicolumn{2}{p{0.3\linewidth}}{
        \cellcolor[rgb]{0.9,0.9,0.9}{
            \makecell[{{p{\linewidth}}}]{
                \texttt{\tiny{[GM$|$GM]}}
                [1, 0, 1]
            }
        }
    }
    & & \\ \\

    \theutterance \stepcounter{utterance}  
    & & & \multicolumn{2}{p{0.3\linewidth}}{
        \cellcolor[rgb]{0.9,0.9,0.9}{
            \makecell[{{p{\linewidth}}}]{
                \texttt{\tiny{[GM$|$GM]}}
                [[1, 0, 1], [0, 1, 1]]
            }
        }
    }
    & & \\ \\

\end{supertabular}
}

\end{document}
